\chapter{언제 탐색하고 언제 결정할까: \nt{NORA}를 더하다}
\chaptermark{\nt{NORA}를 더하다}
\label{sec:nora}

\section{무엇이 빠져 있는가}

\ref{sec:dofa}장에서 네트워크는 잡음 덕분에 학습했다. 활동의 무작위한 편차가 탐색이었고, \nt{DOPA}는 그 편차로 결과가 나아졌는지를 알려 주었다. 그러나 거기서 정하지 않은 매개변수가 하나 남았다. \emph{얼마나} 탐색할 것인가이다. 잡음이 너무 강하면 네트워크는 갈팡질팡하며 이미 아는 것을 활용하지 못한다. 너무 약하면 매번 같은 것만 고르고, 가까운 어딘가가 더 나아졌다는 것을 알아채지 못한다. 강화 학습에서는 이를 활용과 탐색 사이의 선택이라 하며, 모든 알고리즘에는 이를 위한 노브가 있다. \ref{sec:neurontypes}장에서 우리는 뇌에서 이 노브를 돌리는 것이 \nt{NORA}라고 가정했다.

사람의 노르아드레날린 뉴런은 수만 개이고(표~\ref{tab:types}), 거의 모두 작은 무리 하나에 모여 있으며, 그 출력은 사실상 뇌 전체로 퍼진다. 다만 이 무리의 부분마다 서로 다른 영역을 어느 정도 나누어 맡는다는 사실이 밝혀졌다~\cite{Poe2020}. 이것은 \nt{DOPA}보다도 좁은 브로드캐스트 채널이다. 작동 방식은 두 가지다~\cite{AstonJones2005}. \emph{배경} 모드에서 뉴런은 초당 1회 미만에서 몇 회 사이의 발화율로 고르게 스파이크를 내며, 이 발화율은 각성 수준과 함께 천천히 변한다. \emph{돌발} 모드에서는 현재 과제에 중요한 사건에, 대략 결정을 내리는 순간에 짧은 버스트로 응답한다. 편리하지만 정확하지 않은 측정점은 동공 지름이다. 동공은 이 뉴런들의 활동과 함께 변하지만, 다른 영역의 활동~\cite{Joshi2016}과 피질로 가는 콜린성 출력~\cite{Reimer2016}과도 함께 변한다. 동공이 보여 주는 것은 \nt{NORA} 하나가 아니라 전반적인 각성 수준이다.

\section{이득}

측정된 사실은 이것이다. 노르아드레날린은 뉴런의 배경 활동을 자극에 대한 응답보다 더 강하게 억제하여 신호 대 배경 비를 높인다~\cite{Foote1975NE}. 이때 개별 뉴런 특성 곡선의 기울기에 문자 그대로 어떤 계수가 곱해진다는 것까지 실험에서 나오지는 않는다. \ref{sec:neurontypes}장에서처럼 이 효과를 전달하는 단순한 모델을 쓰자. 노르아드레날린은 수신자 전달 특성의 기울기를 바꾼다~\cite{ServanSchreiber1990}. 그러면 약한 문턱값 아래 입력(배경)은 더 작게, 강한 입력은 더 크게 나온다. 뉴런이 다음 발화율로 응답한다고 하자.
\begin{empheq}[box=\keybox]{equation}
r \;=\; F\bigl(g\,(u-\theta)\bigr), \qquad F(z)=\frac{r_{\max}}{1+e^{-z}} ,
\label{eq:gain}
\end{empheq}
여기서 $u$는 입력 전류, $\theta$는 문턱값, $g$는 모델에서 \nt{NORA}가 조절하는 이득이다. $g$가 작으면 발화율은 넓은 범위에서 입력을 매끄럽게 따라간다. 뉴런은 \emph{아날로그 증폭기}다. $g$가 크면 문턱값보다 조금만 높은 입력도 최대 발화율을, 낮은 입력은 0을 낸다. 뉴런은 \emph{비교기}, 거의 디지털 소자다(그림~\ref{fig:gaincurves}). 노브 하나가 같은 뉴런을, 전자공학에서라면 서로 다른 회로가 맡을 두 모드 사이에서 전환한다.

\begin{figure}[t]
\centering
\begin{tikzpicture}[>=Stealth,x=1.1cm,y=2.4cm]
\draw[->,gray!70!black] (-3.3,0) -- (3.4,0) node[below left,black] {\small $u$};
\draw[->,gray!70!black] (-3.3,0) -- (-3.3,1.2) node[left,black] {\small $r$};
\draw[densely dashed,gray!60!black] (0,0) -- (0,1.1);
\node[below,font=\small] at (0,0) {$\theta$};
\draw[densely dotted,gray!60!black] (-3.3,1) -- (3.2,1) node[right,black,font=\small] {$r_{\max}$};
\draw[very thick,blue!60!black] plot[domain=-3.2:3.2,samples=60] (\x,{1/(1+exp(-0.8*\x))});
\draw[very thick,orange!85!black] plot[domain=-3.2:3.2,samples=60] (\x,{1/(1+exp(-2*\x))});
\draw[very thick,red!70!black] plot[domain=-3.2:3.2,samples=120] (\x,{1/(1+exp(min(-8*\x,9)))});
\node[blue!60!black,font=\small,anchor=west] at (0.5,0.12) {작은 $g$: 증폭기};
\node[red!70!black,font=\small,anchor=east] at (-0.3,0.85) {큰 $g$: 비교기};
\end{tikzpicture}
\caption{이득 $g$의 세 값에서의 전달 특성~\eqref{eq:gain}.}
\label{fig:gaincurves}
\end{figure}

큰 이득은 잡음에 도움이 될까? 항상 그렇지는 않다. 두 입력이 $\Delta u$만큼 다르고, 어느 쪽이 큰지 가려야 한다고 하자. 잡음 $\sigma_{\text{pre}}$가 증폭기 \emph{앞}에서 입력에 더해지면 신호와 잡음이 함께 증폭되어 그 비는 변하지 않는다. 반면 잡음 $\sigma_{\text{post}}$가 증폭기 \emph{뒤}에서, 즉 \ref{sec:code}장에서처럼 믿을 수 없는 시냅스와 네트워크 자체의 무작위 스파이크에서 더해지면, 신호 대 잡음비
\begin{empheq}[box=\keybox]{equation}
d' \;=\; \frac{g\,\Delta u}{\sqrt{g^2\sigma_{\text{pre}}^2+\sigma_{\text{post}}^2}}
\label{eq:gainsnr}
\end{empheq}
는 $g\sigma_{\text{pre}}$가 $\sigma_{\text{post}}$와 비슷해질 때까지 $g$와 함께 커진다~\cite{ServanSchreiber1990}.

\begin{keyblock}
\begin{proposition}[이득이 도움이 될 때]
\label{prop:gainnoise}
이득을 높임으로써 \nt{NORA}는 증폭기 \emph{뒤}, 즉 네트워크 안에서 생기는 잡음에 대해서만 입력 구별을 개선한다. 입력과 함께 들어온 잡음은 이득으로 없앨 수 없다.
\end{proposition}
\end{keyblock}

\section{양자택일에서의 이득}

\ref{sec:gaba}장의 선택으로 돌아가자. 입력이 $I_1$, $I_2$인 두 \nt{GLUT} 뉴런 무리가 세기 $\beta$로 서로 억제한다. 이제 각 무리에 이득 $g$가 있다. 선택 방정식~\eqref{eq:wta}에서 괄호 안의 식에 $g$가 곱해진다. 두 무리가 모두 작동하는 동안 정상 상태의 발화율은 $r_1=g(I_1-\beta r_2)$, $r_2=g(I_2-\beta r_1)$에서 구한다. 두 식을 더하고 빼면 합 $r_1+r_2=g(I_1+I_2)/(1+g\beta)$와 차 $r_1-r_2=g(I_1-I_2)/(1-g\beta)$를 얻는다. 둘의 비는 네트워크가 한 입력을 다른 입력보다 얼마나 뚜렷하게 선호하는지를 나타낸다.
\begin{empheq}[box=\keybox]{equation}
\frac{r_1-r_2}{r_1+r_2} \;=\; \varepsilon\,\frac{1+g\beta}{1-g\beta}, \qquad \varepsilon=\frac{I_1-I_2}{I_1+I_2}, \quad g\beta<1 .
\label{eq:wtagain}
\end{empheq}
작은 입력 차 $\varepsilon$은 $(1+g\beta)/(1-g\beta)$배로 증폭되며, 이 인수는 $g\beta$가 1에 다가가면 한없이 커진다. $g\beta>1$이면 명제~\ref{prop:wta}에서처럼 두 무리가 모두 작동하는 상태는 불안정하다. 한쪽이 이기고 다른 쪽은 침묵한다(그림~\ref{fig:pitchfork}).

\begin{figure}[t]
\centering
\begin{tikzpicture}[>=Stealth,x=3.2cm,y=1.5cm]
\draw[->,gray!70!black] (0,0) -- (2.15,0) node[below left,black] {\small $g\beta$};
\draw[->,gray!70!black] (0,-1.25) -- (0,1.35) node[above,black] {\small $(r_1-r_2)/(r_1+r_2)$};
\draw[gray!60!black] (1,-0.04) -- (1,0.04); \node[below right,font=\small] at (1,0) {$1$};
\node[left,font=\scriptsize] at (0,1) {$1$}; \node[left,font=\scriptsize] at (0,-1) {$-1$};
\draw[very thick,blue!60!black] plot[domain=0:0.905,samples=60] (\x,{0.05*(1+\x)/(1-\x)}) -- (1,1) -- (2.05,1);
\draw[very thick,blue!60!black] (1.105,-1) -- (2.05,-1);
\draw[thick,densely dashed,gray!60!black] (1,0) -- (2.05,0);
\node[font=\small,align=center] at (0.45,-0.62) {숙고 중:\\두 무리 모두\\작동};
\node[font=\small,align=center] at (1.55,0.45) {결정함:\\한 무리만 작동};
\node[font=\small,blue!60!black,anchor=south west] at (1.05,-0.97) {약한 입력이 먼저 이겼다: 유지된다};
\end{tikzpicture}
\caption{이득에 따른 양자택일. 입력 차는 $\varepsilon=0.05$이다. $g\beta>1$이면 두 결정이 모두 안정하며(실선; 약한 입력의 승리는 $g\beta>(1+\varepsilon)/(1-\varepsilon)\approx1.1$에서), 네트워크는 이미 내린 결정을 유지한다. 대칭 상태(점선)는 불안정하다.}
\label{fig:pitchfork}
\end{figure}

\begin{keyblock}
\begin{proposition}[이득이 선택 모드를 전환한다]
\label{prop:gainwta}
상호 억제 $\beta$를 가진 선택 네트워크에서 이득 $g$는 네트워크를 경계 $g\beta=1$ 너머로 옮긴다. 경계 아래에서 네트워크는 “숙고”한다. 두 무리가 모두 작동하고 입력 차는 식~\eqref{eq:wtagain}에 따라 증폭된다. 경계 위에서 네트워크는 “결정”했다. 한 무리만 작동하고 결정이 유지된다. \nt{NORA}는 $g$를 바꿈으로써 가중치를 하나도 건드리지 않고 네트워크를 이 두 모드 사이에서 전환할 수 있다.
\end{proposition}
\end{keyblock}

\section{온도로서의 이득}

이제 선택에 네트워크 자체에서, 즉 증폭기 뒤에서 생기는 잡음을 더하자. 어느 무리가 이길지는 $g(u_A-u_B)+\eta$의 부호가 정한다. 여기서 $u_A-u_B$는 입력 차, 즉 \ref{sec:dofa}장에서 학습한 가중치의 차이고, $\eta$는 척도 $\sigma$인 잡음이다. 잡음이 로지스틱 분포를 따른다면(가우스 분포여도 곡선은 거의 같다) $A$를 고를 확률은
\begin{empheq}[box=\keybox]{equation}
P(A) \;=\; \frac{1}{1+e^{-g\,(u_A-u_B)/\sigma}} .
\label{eq:softmax}
\end{empheq}
프로그래머라면 여기서 온도 $T=\sigma/g$인 softmax 선택을 알아볼 것이다. 이득은 역온도다. $g$가 작으면 눈에 띄게 나은 선택지도 못한 선택지보다 그리 자주 뽑히지 않는다. 네트워크는 많이 탐색한다. $g$가 크면 알려진 최선의 선택지가 거의 항상 뽑힌다. 네트워크는 아는 것을 활용한다.

식~\eqref{eq:wtagain}와 \eqref{eq:softmax}의 $g$가 정확히 무엇인지 짚어 두자. 이것은 선택하는 순간 현재 과제 입력 차에 걸리는 \emph{유효} 이득이다. \nt{NORA}의 두 모드는 이 이득에 반대로 작용한다. 결정 순간의 돌발 버스트는 이득을 높여 네트워크가 선택을 굳히게 한다. 반대로 높은 배경 활동은 탐색과 함께 간다. 사람에서는 무작위 선택 직전에 기저 동공이 더 넓고~\cite{JepmaNieuwenhuis2011}, 쥐에서는 앞띠피질로 들어가는 \nt{NORA} 입력이 선택을 무작위 모드로 바꾼다~\cite{Tervo2014stochastic}. 따라서 높은 배경 \nt{NORA}에는 \emph{작은} 유효 $g$가, 버스트에는 큰 유효 $g$가 대응한다.

\section{얼마나 탐색할 것인가}

어떤 이득이 더 나을까? 모델로 확인했다. 네트워크는 식~\eqref{eq:softmax}에 따라 두 행동 가운데 하나를 고른다. 각 행동은 어떤 확률로 보상을 주며, 네트워크는 \ref{sec:dofa}장에서처럼 선택한 행동의 보상으로 그 확률을 추정한다. 때때로 세계가 바뀐다. 두 확률이 새로 정해진다. 바뀌는 것은 선택한 행동일 수도 있고, 네트워크가 오랫동안 시도하지 않은 행동일 수도 있다. 후자는 탐색으로만 알 수 있다. 그림~\ref{fig:invertedU}는 여러 고정 $g$에서 네트워크가 가능한 최선의 보상 중 얼마를 얻는지 보여 준다.

\begin{figure}[t]
\centering
\begin{tikzpicture}[>=Stealth,x=1.2cm,y=1cm]
\draw[->,gray!70!black] (0,0) -- (7.6,0) node[below left,black] {\small $g$};
\draw[->,gray!70!black] (0,0) -- (0,4.4) node[above,black,font=\small] {최선 보상 대비 비율};
\foreach \x/\l in {0/0.5,1/1,2/2,3/4,4/8,5/16,6/32,7/64}{\draw[gray!60!black] (\x,-0.06) -- (\x,0.06); \node[below,font=\scriptsize] at (\x,0) {\l};}
\foreach \v/\y in {0.8/0.8,0.9/2.4,1.0/4.0}{\draw[gray!60!black] (-0.06,\y) -- (0.06,\y); \node[left,font=\scriptsize] at (0,\y) {\v};}
\draw[very thick,blue!60!black] plot[mark=*,mark size=1.3pt] coordinates {(0,0.368) (1,0.880) (2,1.728) (3,2.784) (4,3.456) (5,3.616) (6,3.488) (7,3.264)};
\draw[very thick,red!70!black] plot[mark=*,mark size=1.3pt] coordinates {(0,0.368) (1,0.672) (2,1.184) (3,1.840) (4,2.176) (5,1.920) (6,1.456) (7,1.104)};
\node[blue!60!black,font=\small,anchor=south] at (5,3.7) {조용한 세계};
\node[red!70!black,font=\small,anchor=north] at (5.1,1.25) {빠르게 변하는 세계};
\draw[->,thick,blue!60!black] (5,3.25) -- (5,3.55);
\draw[->,thick,red!70!black] (4,1.8) -- (4,2.1);
\node[font=\small\itshape,gray!35!black,anchor=west] at (0.15,2.3) {맹목적으로 탐색};
\node[font=\small\itshape,gray!35!black] at (6.45,2.55) {변화를 놓침};
\end{tikzpicture}
\caption{고정 이득 $g$에서 가능한 최선의 보상 대비 비율($g$축은 로그 눈금). 파란 곡선은 세계가 평균 1000번 시도에 한 번, 빨간 곡선은 20번에 한 번 바뀌는 경우다. 화살표는 최선의 $g$를 표시한다. 빠르게 변하는 세계에서는 그 값이 절반이다. 4000번 시도의 실행 60회 평균.}
\label{fig:invertedU}
\end{figure}

$g=0.5$에서 네트워크는 지식을 거의 활용하지 못해 가능한 보상의 약 4분의 3을 얻고, $g$가 매우 크면 변화를 놓친다. 최선의 값은 세계가 1000번 시도에 한 번 바뀔 때 $g=16$(비율 $0.976$), 20번에 한 번 바뀔 때 $g=8$($0.886$)이다.

\begin{keyblock}
\begin{proposition}[뒤집힌 U]
\label{prop:explore}
네트워크가 모으는 보상은 이득에 대해 뒤집힌 U 모양으로 의존한다. $g$가 너무 작으면 시도를 탐색에 낭비하고, 너무 크면 변화를 알아채지 못한다. 세계가 빨리 변할수록 최선의 $g$는 작다.
\end{proposition}
\end{keyblock}

뒤집힌 U는 실험에서도 관찰된다. 동물은 \nt{NORA} 뉴런의 배경 활동이 적당하고 돌발 응답이 뚜렷할 때 과제를 가장 잘 풀고, 배경 활동이 높으면 산만해져 과제 사이를 오간다~\cite{AstonJones2005}. 이는 앞 절에서 본 두 모드의 차이와 들어맞는다. 결정 순간의 돌발 버스트는 바로 선택해야 할 때 큰 유효 $g$를 준다. 버스트 없는 높은 배경 활동은 무관한 입력까지 모든 것을 증폭하므로 현재 과제에 대한 유효 $g$가 떨어진다. 이것이 탐색이다.

\section{누가 노브를 돌리는가}

최선의 이득이 세계가 얼마나 빨리 변하는지에 달려 있다면 이득을 맞추어 조정하는 편이 좋다. 그런데 \nt{NORA} 뉴런은 세계가 바뀌었다는 것을 어떻게 알까? 자연스러운 표지는 \emph{뜻밖의 일}이다. 예측 오차가 평소보다 커지는 것이다. 여기서 두 종류의 불확실성을 구별해야 한다. 보상이 그냥 무작위라면 오차는 늘 크고, 그것은 예상된 일이다. 반면 오차가 평소 크기보다 갑자기 커졌다면 무언가가 바뀐 것이다. 한 가설에 따르면 \nt{NORA}가 알리는 것은 바로 이 예상치 못한 불확실성이고, 예상된 불확실성은 \nt{ACET}가 알린다~\cite{YuDayan2005}.

이 신호로 무엇을 할까? 이득을 낮추어 더 탐색할 수 있다. 학습률을 높여 낡은 추정을 더 빨리 잊을 수도 있다. 실험에 따르면 급격한 변화 뒤에 사람은 더 빨리 학습하고, 이때 동공이 커진다~\cite{Nassar2012}. 다시 말하지만 동공은 \nt{NORA}만이 아니라 \nt{ACET}의 지표이기도 하다~\cite{Reimer2016}. 둘을 한꺼번에 할 수도 있다. 또 다른 가설에 따르면 \nt{NORA}의 돌발 버스트는 \emph{인터럽트}로 작동한다. 네트워크가 현재 상태를 초기화하고 새 상황에 맞게 재구성하게 하는 신호로~\cite{BouretSara2005}, 프로세서의 공통 인터럽트 선과 같다.

찾지 못한 것도 솔직히 말해 두자. 모델에서 우리는 두 가지 단순한 조정 방식을 모두 시험했다. 평활한 오차가 장기 평균을 넘으면 $g$를 낮추거나 학습률을 높이는 것이다. 한동안 오래 가만히 있다가 한동안 자주 바뀌는 세계에서, 두 방식의 최적 설정은 최선 보상의 $0.926$--$0.927$을 얻었다. 최선의 고정 $g$($g=8$에서 $0.925$)나 고정이지만 충분히 큰 학습률($0.928$)과 거의 같고, 설정이 나쁘면 그보다 적었다. 그림~\ref{fig:invertedU}의 곡선은 꼭대기 근처에서 완만하므로, 이런 세계에서는 조정할 것이 거의 없다. 조정이 값을 하려면 모드마다 전혀 다른 설정이 필요해야 한다. \nt{NORA}가 어떤 제어 법칙을 구현하는지는 아직 밝혀지지 않았다.

\section{정리}

\begin{table}[t]
\centering
\footnotesize
\setlength{\tabcolsep}{4pt}
\renewcommand{\arraystretch}{1.15}
\begin{tabular}{>{\raggedright}p{3.0cm}>{\raggedright}p{4.4cm}>{\raggedright\arraybackslash}p{6.0cm}}
\toprule
\nt{NORA}가 하는 일 & 방법 & 알려진 것 \\
\midrule
뉴런 응답의 기울기를 바꾼다 & 식~\eqref{eq:gain}의 이득 $g$ & 신호 대 잡음비 증가는 측정됨; 곱셈 모델은 단순화 \\
선택을 전환한다 & 네트워크를 $g\beta=1$ 너머로 옮긴다 & 모델~\eqref{eq:wtagain}에서 나옴; 문턱값의 직접 측정은 없음 \\
얼마나 탐색할지 정한다 & $g$는 역온도~\eqref{eq:softmax}; 배경은 작은 유효 $g$(탐색), 버스트는 큰 $g$(결정) & 뒤집힌 U와 두 작동 모드는 측정됨; 탐색과의 연관은 근거가 탄탄한 가설 \\
변화를 알린다 & 예상치 못한 불확실성 & 변화 뒤 동공과 학습률 증가는 측정됨; 그것이 \nt{NORA} 신호라는 것은 가설 \\
끼어들어 초기화한다 & 뜻밖의 사건에 대한 돌발 버스트 & 중요한 사건에 대한 버스트는 측정됨; “인터럽트”는 가설 \\
\bottomrule
\end{tabular}
\caption{\nt{GLUT}, \nt{GABA}, \nt{DOPA}로 된 네트워크에 \nt{NORA}가 더하는 것.}
\label{tab:nora}
\end{table}

\nt{DOPA}는 네트워크에 가중치를 \emph{어느 쪽으로} 바꿀지 알려 준다. \nt{NORA}는 가중치를 하나도 바꾸지 않지만, 네트워크가 이미 아는 것을 \emph{어떻게} 쓸지 정한다(표~\ref{tab:nora}). 노브 하나, 즉 이득이 뉴런을 증폭기에서 비교기로, 선택 네트워크를 “숙고 중”에서 “결정함”으로, 선택을 탐색에서 활용으로 바꾼다. \nt{NORA}가 세계의 변화 속도를 어떻게 아는지는 열린 문제다.

\section{연습문제}

\begin{exercise}[\lvlA]
\label{ex:08:1}
\emph{뇌 전체에 노브 하나.} (a)~표~\ref{tab:types}로 \nt{NORA} 뉴런 하나당 뇌 뉴런이 몇 개인지 추정하라. (b)~증폭기 앞 잡음은 $\sigma_{\text{pre}}=0.5$, 뒤 잡음은 $\sigma_{\text{post}}=4$이다. 식~\eqref{eq:gainsnr}의 $d'$이 극한의 $90\,\%$에 이르는 $g$는 얼마인가?
\end{exercise}

\begin{exercise}[\lvlB]
\label{ex:08:2}
\emph{이득이 있는 선택.} $\tau\,\dd r_1/\dd t=-r_1+g[I_1-\beta r_2]_+$, $\tau\,\dd r_2/\dd t=-r_2+g[I_2-\beta r_1]_+$, $\tau=20\ms$. (a)~$g\beta=0.5$와 $0.9$에서 $r_1-r_2$가 감쇠하는 시간은? (b)~$\varepsilon=0.05$에서 두 무리가 모두 작동하는 $g\beta$의 상한과, 약한 입력의 승리가 안정해지는 $g\beta$의 하한을 구하라.
\end{exercise}

\begin{exercise}[\lvlA]
\label{ex:08:3}
\emph{잡음에서 나오는 softmax.} $K$개 무리가 저마다 굼벨 잡음 $P(\eta_k<z)=\exp(-e^{-z/\sigma})$를 받고, $gu_k+\eta_k$가 가장 큰 무리가 이긴다. $P(k)$를 구하라. $u_A-u_B=0.1$, $\sigma=1$인 두 선택지 중 나은 쪽이 $90\,\%$ 선택되는 $g$는?
\end{exercise}

\begin{exercise}[\lvlB]
\label{ex:08:4}
\emph{최선 이득의 추정.} 그림~\ref{fig:invertedU}의 모델에서 변화 뒤 보상 확률은 $[0,1]$에서 균등하다. 네트워크는 $e^{g\Delta}$번에 한 번 못한 행동을 시도한다. 이를 $1/h$와 같게 두면 $g^*\approx\ln(1/h)/\bar\Delta$이다. $\bar\Delta=\langle|p_A-p_B|\rangle$와 $h=0.001$, $0.01$, $0.05$에서의 $g^*$를 구하라.
\end{exercise}

\begin{exercise}[\lvlB]
\label{ex:08:5}
\emph{시뮬레이션: 이득과 변화 속도.} \texttt{models/\allowbreak nora\_explore.py}의 사본(현재 폴더에 파일을 쓴다)으로 $h=0.001$부터 $0.2$까지 $g^*(h)$를 구하라. 연습문제~\ref{ex:08:4}의 추정은 어디서 맞는가? 변화가 빠를 때 $g^*$가 더 줄지 않는 이유는?
\end{exercise}

\begin{exercise}[\lvlB]
\label{ex:08:6}
\emph{설계: 선택 네트워크를 위한 인터럽트.} 연습문제~\ref{ex:08:2}의 네트워크($\beta=2$, $\tau=20\ms$, $g=1$)는 입력이 이제 $I_1=0.9$, $I_2=1.0$인데도 $r_1=0.9$, $r_2=0$을 유지한다. \nt{NORA}가 시간 $T_p$ 동안 $g$를 $g_{\text{lo}}$로 낮춘다. 최소 $T_p$를 $g_{\text{lo}}=0$에서는 해석적으로, $g_{\text{lo}}=0.25$에서는 시뮬레이션으로 구하라.
\end{exercise}

\begin{exercise}[\lvlC]
\label{ex:08:7}
\emph{조정이 값을 할 때.} 오라클은 지금이 조용한 시기인지 변덕스러운 시기인지 알고 시기마다 $g$를 따로 정한다. (a)~이 장의 세계(시기당 $1000$번 시도, $h=0.001$과 $0.05$)에서 최선의 고정 $g$ 대비 이득을 시뮬레이션으로 구하라. (b)~이득이 더 큰 세계를 만들고, 뜻밖의 일에 따른 조정이 그중 얼마를 가져가는지 구하라.
\end{exercise}
